\chapter{Conclusion}

Ce dernier chapitre répond directement à la question posée en
introduction, puis donne un regard d'ensemble sur ce que la chaîne Big
Data a produit. Les constats détaillés et les recommandations ne sont
pas repris ici : l'objectif est de prendre du recul sur l'ensemble, pas
de refaire la synthèse.

\section{Réponse à la question directrice}

La question était : \emph{quelles tendances du marché des emplois IA
peut-on mettre en évidence à partir des offres, des métiers, des
compétences et des rémunérations ?} Après avoir transformé 1500 offres
brutes en un modèle exploitable de 25 métiers et 9\,428 liaisons
offre-compétence, notre réponse tient en une idée : sur ce dataset, le
marché de l'emploi IA se distingue moins par des écarts salariaux
spectaculaires entre profils que par la \textbf{concentration de la
valeur sur quelques signaux précis}. Un métier rare et massivement
demandé (LLM Engineer), une compétence universelle (Python), et un mode
de travail devenu flexible par défaut (hybride ou remote) structurent le
marché bien plus fortement que la seule ancienneté.

Quelques chiffres pour fixer les ordres de grandeur :

\begin{center}
\begin{tabular}{ll}
\toprule
Indicateur & Valeur (dataset complet) \\
\midrule
Offres analysées & 1500 \\
Métiers distincts & 25 \\
Salaire moyen / médian & 194\,892\$ / 179\,421\$ \\
Compétence la plus demandée & Python (942 offres, 62,8\,\%) \\
Métier le plus tendu & LLM Engineer (75 offres, 240\,960\$ en moyenne) \\
Pays le plus représenté et le mieux payé & USA (515 offres, 226\,190\$) \\
Mode de travail dominant & Hybrid (45,7\,\%) \\
Rôles LLM/GenAI & 327 offres (21,8\,\%) \\
Écart de salaire au sein du top 10 pays & 20,2\,\% seulement \\
\bottomrule
\end{tabular}
\end{center}

Ce que ces chiffres racontent, une fois reliés, c'est un marché où la
rémunération traduit surtout la tension d'une compétence ou d'un profil,
et assez peu l'accumulation d'années d'expérience (corrélation de 0,079
contre 0,292 pour la demande). Le salaire reste élevé partout dans le
haut du tableau, sans abîme entre les pays : le marché mondial de l'IA
semble se jouer davantage sur la capacité à attirer les bons profils que
sur l'arbitrage géographique.

\section{Apports du projet}

Au-delà des résultats métier, le projet a produit une chaîne complète et
reproductible : ingestion HDFS, tables Hive RAW et processed,
transformations Spark, requêtes Hive SQL et Spark SQL, analyses KPI,
visualisations Jupyter/PySpark. Chaque étape est chiffrée, versionnée et
testée. Le passage systématique par HDFS et Hive garantit que les
analyses ne s'appuient jamais sur une copie locale du CSV, et que les
données brutes restent immuables.

\section{Perspectives}

À titre de suite, une comparaison systématique à des sources externes
(benchmarks salariaux, offres d'emploi temps réel) permettrait de juger
de la représentativité du dataset. Une actualisation régulière des
offres et l'intégration de séries temporelles plus longues seraient
également nécessaires pour transformer les constats descriptifs en
indicateurs de tendance fiables.

La réponse à la question directrice est donc nuancée : oui, des
tendances nettes existent (concentration de la demande, prééminence des
compétences sur l'ancienneté, flexibilité par défaut), mais elles valent
pour \emph{ce} dataset, à \emph{cet} instant. Le principal apport du
projet est une chaîne reproductible, testée et chiffrée, capable de
produire ces constats et de les réviser si la donnée change.
